\section{Functional inputs at the unchanged control}\label{sec:functional-hardness}

The standard control can already support an undecidable comparison problem
when free functional variables are admitted.  This is a different domain
from the all-ground question.  We give the reduction because the distinction
identifies exactly where its additional arithmetic information enters.

The published undecidable problem used here is solvability of a finite
system of linear equations
\[
 \sum_i a_{ji}(X)p_i(X)-g_j(X)=0,
 \qquad a_{ji},g_j\in\mathbb Z[X],\quad p_i\in\mathbb N[X].
\]
It is linear in the unknown polynomials, not a system of arbitrary
Diophantine products of unknowns.  Narendran's theorem and Dong's discussion
give this precise polynomial-semiring boundary \cite{narendran,dong}.
Zero polynomials are allowed.  We use undecidability, without importing a
stronger many-one completeness classification of that source problem.

Put $\tau=o\to o$.  For actual sources $A,B:\tau$ define
\begin{align*}
 \mathsf P(A,B)&=(\lambda z:\tau.A)B,\qquad
 \mathsf H(A)=R^o0(\lambda n:o.A),\\
 \mathsf S(A,B)&=(\lambda z:(o\to\tau).A)(\lambda n:o.B).
\end{align*}
The original argument order is
$R^\rho:o\to(o\to\rho\to\rho)\to\rho\to\rho$:
its first argument $0:o$ is the recursion test, not its base value.
Thus $R^\rho0(\lambda n:o.A)$ has type $\rho\to\rho$
when $A:\rho\to\rho$, also at the higher type used below.
The displayed bound names are unused and fresh.  Writing
$a_i=\can A_i$, $b_i=\can B_i$, the literal canonical clauses give
\begin{align*}
 \can{\mathsf P(A,B)}_1&=8a_1+4b_1+8,\\
 \can{\mathsf S(A,B)}_1&=8a_1+8b_1+16,\\
 \can{\mathsf S(A,B)}_0
   &=\psi\bigl(\omega\otimes(8a_1+8b_1+16)+9a_0+5b_0+13\bigr),\\
 \can{\mathsf H(A)}_1&=2\omega\otimes a_1+3\omega.
\end{align*}
For example the unused $\tau$ binder has vector
$(9a_0+8,2a_1+2,2)$, and its box with $B$ gives the first formula.
For the second binder its scalar is again $9a_0+8$ and its first field
$2a_1+2$; the ground lambda on $B$ has vector
$(5b_0+4,2b_1+2)$, giving both $\mathsf S$ formulas.
The recursor prime followed by its two displayed arguments gives
the $\mathsf H$ formula.  These calculations use actual source syntax;
they do not introduce an arbitrary ordinal-valued primitive.

\begin{lemma}[Signed polynomial pairs]\label{lem:pair-compiler}
Every signed linear expression $D$ in ordinal unknowns with coefficients
in $\mathbb Z[\omega]$ has effectively constructed actual $A_D,B_D:\tau$
whose first-coordinate difference is $2^kD$ for a finite $k$.
All free functional variables have type $\tau$, and internal types have
level at most two.  Differences are taken only in the signed Cantor-monomial
group; each of the two source coordinates remains a positive ordinal.
\end{lemma}
\begin{proof}
Let $C=\lambda u:o.u$, with $C_1=1$.  The first fields of
$\mathsf P(C,C)$ and $\mathsf S(C,C)$ are 20 and 32.  Take sixteen pairs:
eight copies of $(\mathsf S(C,C),\mathsf P(C,C))$ of difference 12,
five of $(C,\mathsf P(C,C))$ of difference $-19$, and three of $(C,C)$.
Combine them componentwise in a balanced $\mathsf S$ tree of height four.
Constants cancel in every paired combination, so the resulting difference is
\[
 8^4(8\cdot12-5\cdot19)=2^{12}.
\]
For a free $X:\tau$, the pair
$(\mathsf P(X,C),\mathsf P(C,C))$ has difference $8(x_1-1)$.
Three multiplications by eight followed by a paired $\mathsf S$ combination
with the constant pair give $2^{15}x_1$.

The pair operation $(A,B)\mapsto(\mathsf S(A,C),\mathsf S(B,C))$
multiplies a difference by eight, and
$(A,B)\mapsto(\mathsf S(A,A),\mathsf S(B,B))$ multiplies it by sixteen.
Applying $\mathsf H$ to both entries multiplies a difference by $2\omega$;
swapping entries negates it.  Repeated addition constructs every finite
integer coefficient.  To add pairs with scale exponents $k_1,k_2$, choose
$K\ge\max(k_1,k_2)+6$.  Every integer at least six is $3u+4v$ with
$u,v\ge0$, so the two scale operations align both differences at $2^K$.
A paired $\mathsf S$ then gives their sum with scale $2^{K+3}$.
Induction on the finitely many signed monomials proves the assertion.
The displayed source combinators introduce no internal level above two.
\end{proof}

The two scalar coordinates of an equation pair need not agree when their
positive coordinates agree.  We therefore pay that alignment explicitly.

\begin{lemma}[Alignment and a full-vector detector]\label{lem:functional-detector}
The condition $\can A_1=\can B_1$ is equivalent to existence of two fresh
positive ground inputs making the full canonical vectors of
\[
 A'=(\lambda t:o.A)z,\qquad B'=(\lambda t:o.B)w
\]
equal.  Full-vector equality of any actual pair of type $\tau$ can in turn
be detected by ground sources $L,R$ satisfying $L_0\ge R_0$, with equality
exactly at vector equality and otherwise gap at least two.
\end{lemma}
\begin{proof}
The aligned first fields are $2a_1+2$ and $2b_1+2$, and their scalars are
\[
 \psi(\omega\otimes(2a_1+2)+5a_0+4+z_0),\qquad
 \psi(\omega\otimes(2b_1+2)+5b_0+4+w_0).
\]
Equal vectors imply $a_1=b_1$.  If this equality holds, choose
$z_0=5b_0+1$, $w_0=5a_0+1$.  Both tails and both vectors then agree.
These choices are made after the original valuation, so are not circular.

For the detector, now write $A,B$ for the pair being tested, and form
\[
 X=\mathsf S(A,A),\quad Y=\mathsf S(B,B),\quad
 Z=\mathsf S(A,B),\quad W=\mathsf S(B,A).
\]
Make the actual shared ground payload
$Q=W(Z(Y(X0)))$, and put $q=Q_0$.
The unconditional finite-tail lower bound for $\psi$ gives
\[
 q\ge1+x_0+y_0+z_0+w_0\ge5.
\]
Set $g_X=XQ$, and similarly $g_Y,g_Z,g_W$, and finally set
$L=e(g_X,g_Y)$, $R=e(g_Z,g_W)$.

Suppose first $a_1>b_1$.  Then $x_1>z_1=w_1$.  The ordinals
\[
 \alpha=\omega\otimes z_1+z_0+q+2,
 \qquad \beta=\omega\otimes x_1+x_0+q
\]
satisfy $\alpha<\beta$.  We must check a norm premise before comparing
their collapses.  Bounded norms of actual assignments and the published
natural-product norm inequality \cite[Propositions 2.2--2.3]{ww} give
\[
 \norm(\alpha)\le2\norm(z_1)+z_0+q+2
 \le3z_0+q+2\le4q+2
 \le2^{q+1}\le\Phi(\norm(\beta)).
\]
Controlled strictness therefore gives $\psi(\beta)>\psi(\alpha)$.
Two consecutive fixed-limit tail-growth steps from \cite{wwfoundations}
yield $g_{X,0}>4g_{Z,0}$.  The same calculation with $W$ gives
$g_{X,0}>4g_{W,0}$.  If $b_1>a_1$, use $Y$ in place of $X$.

If $a_1=b_1$ but $a_0\ne b_0$, the four $\mathsf S$ heads coincide.
Their scalar tails are
\[
 14a_0+13,\quad14b_0+13,\quad
 9a_0+5b_0+13,\quad9b_0+5a_0+13.
\]
The larger diagonal tail exceeds both cross tails by at least five.
Strict increase gives at least a two-unit scalar gap, so two fixed-tail
growth steps in the subsequent applications give again a diagonal $g$
value greater than four times each cross value.

In either unequal case,
$3g_{X,0}+g_{Y,0}>3g_{Z,0}+g_{W,0}$: even when the larger diagonal
is the one with coefficient one, its fourfold domination suffices.
The final eraser is increasing and has gaps at least two.
Equal input vectors give four equal $\mathsf S$ vectors, hence equality.
This proves the claimed detector on every admissible valuation, not only
on the specially chosen alignment witnesses.
\end{proof}

\begin{theorem}[Standard-control functional-input undecidability]
\label{thm:functional-hardness}
Full universal canonical weak comparison and strict comparison are
undecidable already for ground output, free-variable levels at most one,
and internal levels at most two.  The reduction at this level uses finitely
many comparison queries.  A computable single-query reduction from the
same polynomial-system unsatisfiability problem is available with internal
level at most three.
\end{theorem}
\begin{proof}
Let $\pi$ retain exactly the Cantor monomials with finite exponents.
It maps ordinals below $\varepsilon_0$ to $\mathbb N[\omega]$.
It preserves natural sum and multiplication: a natural sum of exponents
is finite exactly when both exponents are finite.  It consequently extends
to the signed Cantor-monomial group and respects multiplication by fixed
$\mathbb Z[\omega]$ coefficients.

For the level-two assertion, enumerate the finitely many possible sets of
zero polynomial unknowns, substitute zero for those, and require the others
to be positive ordinal first coordinates.  Compile the resulting equations
by Lemma~\ref{lem:pair-compiler}, align them separately, and apply
Lemma~\ref{lem:functional-detector}.  Aggregate the nonnegative ground
detectors in identically shaped eraser trees.  Root equality is exactly
simultaneous equation satisfaction.  Its absence is one universal strict
comparison; the $KR\mle IL$ conversion in Theorem~\ref{thm:compiler}
gives the same weak query.  An empty system receives a tautology.

A polynomial solution lifts in the matching zero pattern: each remaining
nonzero polynomial is an admissible positive first field, after choosing its
zeroth field sufficiently large for its norm.  Conversely any positive-ordinal
solution projects by $\pi$ to a polynomial solution of the original system.
Some projected unknowns may become zero, which is permitted in that system.
Thus an algorithm for either comparison would decide polynomial solvability
using the finitely many patterns, contradicting the published theorem.

For a single query choose instead a fixed positive ordinal $\Omega$ with
$\pi(\Omega)=0$, and transform each original equation to
\[
 \sum_i a_i(\omega)x_i-g(\omega)
                    -\left(\sum_i a_i(\omega)\right)\Omega=0.
\]
A polynomial solution $u_i$ lifts by $x_i=u_i\oplus\Omega>0$.
Conversely projection of any positive-ordinal solution solves the original
equations.  The transformation preserves simultaneous use of each unknown.

We must realize the extra constant in actual source syntax.  Put
$\sigma=\tau\to\tau$, $C_2=\lambda u:\tau.u$, and
$\mathsf H_2(A)=R^\tau0(\lambda n:o.A)$.  The canonical clauses give
$(C_2)_2=2$ and $(\mathsf H_2(A))_2=\omega+2\omega\otimes A_2$.
The closed $H=\mathsf H_2(\mathsf H_2(C_2))$ has
\[
 H_2=10\omega^2+\omega,
 \qquad\Omega=2^{10\omega^2+\omega}=\omega^{10\omega+1},
 \qquad\pi(\Omega)=0.
\]
Apply this same actual $H$ to the two entries of the constant pair of
Lemma~\ref{lem:pair-compiler}.  The box gives
$\can{HA}_1=\Omega\otimes(H_1+A_1)$, so their difference is
$2^{12}\Omega$.  Add this as a second closed seed to the pair compiler.
Only fixed $\omega$ shifts and signed sums of that seed are needed.
Compile the one transformed system and apply the same alignment and detector.
The closed $R^\tau$ raises the internal level bound to three, while every
free variable still has level at most one.  This is the promised single-query
reduction.  No claim about all-ground free variables follows.
\end{proof}

\section{An additive-one change of control}\label{sec:control-hardness}

Growth alone does not explain the positive standard-control results in
\cite{wwmultivariate}.  Their modular hypotheses carry indispensable
information.  The following construction changes the control while retaining
its growth, any prescribed finite prefix, and even pointwise additive accuracy.

Fix a primitive recursive bounded simulation predicate $H(e,t)$ for an
effective enumeration of blank-input programs: it is one precisely when
program $e$ has halted within $t$ steps.  For a fixed finite $b$, define
\[
 \epsilon_b(0)=0,\qquad
 \epsilon_b(n)=
 \begin{cases}0&0<n\le b,\\H(v_2(n),n)&n>b,\end{cases}
 \qquad \Gamma_b(n)=\Phi(n)+\epsilon_b(n),
\]
where $v_2(n)$ is the exponent of two dividing $n>0$.
The one function $\Gamma_b$ is fixed before a program index is supplied.

\begin{lemma}[Control properties]\label{lem:perturbed-control}
The function $\Gamma_b$ is primitive recursive, agrees with $\Phi$ on
$[0,b]$, and satisfies
\[
 0\le\Gamma_b(n)-\Phi(n)\le1,\quad
 \Phi(n)\le\Gamma_b(n)<\Phi(n+1),\quad
 \Gamma_b(n+1)>4\Gamma_b(n).
\]
It has the additive margin
$\Gamma_b(v+a)\ge\Gamma_b(v)+2a$ for $a\ge1$, is eventually bounded
by $F_6$, and $\Gamma_b(n)/\Phi(n)\to1$.
\end{lemma}
\begin{proof}
Bounded simulation, valuation $v_2$, and each fixed $F_j$ are primitive
recursive.  The prefix and additive estimates follow from the definition.
If $u=\Phi(n)$, the control estimate in \cite{wwfoundations} gives
$\Gamma_b(n+1)\ge2^{2^u}>4(u+1)\ge4\Gamma_b(n)$.
It also gives $\Phi(n)+1<\Phi(n+1)$ and increments greater than two.
Sum those increments for the additive margin.  For $n\ge101$,
$F_5(n+101)\le F_5(F_5(n))\le F_6(n)$, since $F_5(n)\ge n+101$;
this bounds $\Gamma_b(n)<F_5(n+101)$.  The bounded error and
$\Phi(n)\to\infty$ give the ratio limit.
\end{proof}

These estimates supply the identity-majorization and additive margin of the
collapse construction, as well as its eventual finite-hierarchy upper bound.
Only the coefficient-one collapse will be used for hardness.  For a strictly
increasing identity-majorizing control its exact formula is
$\psi_{\Gamma_b}(\omega+n)=\Gamma_b(n+2)+1$, by the maximum-predecessor
recurrence.  It is injective.  Its consecutive values exceed four times the
preceding ones at the arguments used by the eraser.  Consequently the source
addition, constants, collision detector and conjunction construction proved
above work verbatim with this control.  We do not transfer a standard-control
residue theorem to $\Gamma_b$.

\begin{theorem}[Ground hardness within distance one]\label{thm:control-hardness}
For every $b$, universal canonical weak comparison and strict comparison
at the one fixed control $\Gamma_b$ are $\Pi^0_1$-complete already on
ground-output sources having only ground free names and ground binders,
using only the prime constants $0,S$, with internal levels at most one.
Moreover $\Gamma_b(n)\bmod2$ is not eventually periodic.
\end{theorem}
\begin{proof}
Every value of $F_j$ at a positive argument, and thus every value of
$\Phi$, is a power of two divisible by two.  This follows by induction
from its base and its positive number of iterations.
For $n>b$ the value $\Gamma_b(n)$ is therefore odd exactly when
$H(v_2(n),n)=1$.

Fix the odd constant $C=2b+5$.  For a program index $e$ consider the
existential conjunction over natural variables
\[
 z+4=2^e(2t+C),\qquad y=F_S(z),\qquad y=w+w+1.
\]
Here $F_S(z)=\Gamma_b(z+4)$, and $2^e$ is a fixed coefficient computed
from the input index.  It is not a variable exponentiation operation.
With $n=2^e(2t+C)$ one has $n\ge5$, $n>b$ and $v_2(n)=e$.
The conjunction has a solution at $t$ exactly when $H(e,n)=1$.
If $e$ halts, choose $t$ large enough that $n$ exceeds its halting time;
then $z=n-4\ge0$ and $w=(\Gamma_b(n)-1)/2\ge0$ supply witnesses.
Conversely any solution certifies the bounded halting test.  Thus the
conjunction is satisfiable exactly when program $e$ halts.

Flatten it using the proved addition, constant and actual $S$ graph gadgets,
then use the collision detector and eraser aggregation.  Its unsatisfiability
is one universal strict comparison and, by $KR\mle IL$, one weak comparison.
All introduced binders are ground, and all leaves use only $0,S$ and ground
variables.  The nonhalting reduction proves $\Pi^0_1$ hardness with the stated
source restrictions.  The arithmetic DAG uses $O(e+\log(b+2))$ gates;
this is a bound in the numerical index $e$, not its binary length.

For membership, enumerate all finite tuples of positive raw ground values.
Every value of a displayed source is computable: its closed heads are finite
and the collapse reduces to fixed finite iteration of the primitive recursive
control.  A failed universal comparison is therefore semidecidable.
This proves completeness, without an oracle hidden in $\Gamma_b$.

Finally suppose parity were periodic with period $p\ge1$ after $N$.
For each $e$ effectively choose $n_e>\max(N,b)$ with $v_2(n_e)=e$.
Every point
\[
 n_e+k p2^{e+1},\qquad k=0,1,2,\ldots,
\]
has the same residue modulo $p$ and the same valuation $e$.
Periodicity would make the bounded tests $H(e,n_e+kp2^{e+1})$ constant.
They are always zero if $e$ never halts and eventually one if it halts.
Computing $\Gamma_b(n_e)\bmod2$ would thus decide halting.
The alleged finite $N,p$ can be hardwired; they need not be found uniformly
from $b$ for this contradiction.  No eventual parity period exists.
\end{proof}

The functional-input theorem leaves the control unchanged and changes the
input domain.  The last theorem keeps ground inputs but changes the control.
Neither is a negative answer for the remaining exact-standard-control,
all-ground canonical comparison problem.  Together with the positive
standard-control classes, they explain why a claimed solution must preserve
both the source domain and the precise arithmetic function.
